% 제4장: 실험 결과 및 분석 / 4.2 문체 수정본
% section4_1.tex 다음에 삽입한다. 표는 4.4에서도 함께 참조한다.
% 근거: results/koblex-reproduction-v1/analysis.json,
%       results/koblex-reproduction-v1/chapter4-prep-v1/A1/numeric-ledger.json,
%       results/koblex-reproduction-v1/chapter4-prep-v1/A4/stage-summary.jsonl,
%       results/koblex-reproduction-v1/chapter4-prep-v1/A4/pool-feasibility-summary.json

\subsection{재정렬 모델 교체에 따른 근거 선택 성능}
\label{sec:ch4_evidence_results}

Jev의 최종 근거 F1은 61.75점으로 수정 CE의 58.57점보다 3.19점
높았다(표~\ref{tab:ch4_main_results}). 차이의 탐색적 95\% 신뢰구간은
$[+0.75,\,+5.68]$점이었다. Recall도 57.74점에서 60.69점으로
2.95점 증가하였으며, 차이의 신뢰구간은 $[+0.39,\,+5.51]$점이었다.
두 지표는 Jev 조건에서 정답 근거의 평균 회수 수준이 높았음을 보여준다.

근거 선택 점수는 공식 정답 근거(gold)와 최종 선택 근거의 문자열
집합을 비교한 문항별 점수의 비가중 평균이다. EM은 두 집합의 완전
일치율이며, 모든 gold 포함률은 추가 근거의 유무와 관계없이 정답
근거를 모두 포함한 문항의 비율이다.

\begin{table}[htbp]
    \centering
    \caption{수정 CE와 Jev의 전체 성능 비교 ($n=226$)}
    \label{tab:ch4_main_results}
    \small
    \setlength{\tabcolsep}{3pt}
    \begin{tabular}{@{}lrrrr@{}}
        \hline
        지표 & 수정 CE & Jev & $\Delta$ & 차이의 95\% CI \\
        \hline
        \multicolumn{5}{@{}l}{\textbf{근거 선택}} \\
        F1              & 58.57 & 61.75 & $+3.19$ & $[+0.75,\,+5.68]$ \\
        Recall          & 57.74 & 60.69 & $+2.95$ & $[+0.39,\,+5.51]$ \\
        EM              & 27.88 & 29.65 & $+1.77$ & $[-1.77,\,+5.19]$ \\
        모든 gold 포함률 & 34.07 & 34.96 & $+0.88$ & $[-2.90,\,+4.53]$ \\
        \hline
        \multicolumn{5}{@{}l}{\textbf{답변 품질}} \\
        Token F1        & 46.20 & 47.31 & $+1.12$ & $[+0.02,\,+2.27]$ \\
        LF-Eval         & 67.57 & 70.15 & $+2.58$ & $[+0.37,\,+4.79]$ \\
        \hline
    \end{tabular}
    \par\smallskip
    \begin{minipage}{\linewidth}
        \footnotesize
        값은 0--1 원점수의 100배이며, EM과 모든 gold 포함률은 백분율이다.
        $\Delta$는 Jev$-$수정 CE로, EM과 모든 gold 포함률의 차이는 \%p,
        나머지 지표의 차이는 점 단위이다. 차이와 CI는 반올림 전 값으로
        계산하였다. CI는 151개 의존 그룹의 대응 bootstrap 10,000회로
        산출한 탐색적 percentile 구간이며, 다중 비교를 보정하지 않았다.
    \end{minipage}
\end{table}

정답 근거를 모두 확보하는 성능의 개선은 불확실했다.
EM의 증가폭은 1.77\%p였고 신뢰구간은 0을 포함하였다.
모든 gold를 포함한 문항은 수정 CE 77개에서 Jev 79개로 늘었으나,
비율 차이는 0.88\%p에 그쳤으며 신뢰구간은
$[-2.90,\,+4.53]$\%p였다. 평균 F1과 Recall의 개선에 비해
근거 집합의 완결성에서는 뚜렷한 차이를 확인하기 어려웠다.

Jev의 회수 이점은 LLM 선택기에 제시하는 shortlist에서도 관찰되었다.
공통 BM25 후보 합집합의 문항 평균 gold Recall은 77.21\%,
모든 gold 포함률은 61.06\%(138/226)였다.
예상 법령별 상위 10개를 합친 shortlist의 평균 gold Recall은
수정 CE 69.47\%, Jev 74.63\%로, Jev가 더 많은 정답 근거를
후속 선택 단계에 남겼다. 후보 목록당 하나를 선택하는 제약까지
고려할 때, 모든 gold를 동시에 선택할 수 있는 문항은 수정 CE 104개,
Jev 120개였다. 이는 구조적으로 가능한 문항 수의 상한이다.
실제 완전 회수 문항은 각각 77개와 79개로 이보다 적었다.
따라서 최종 성능은 변경된 재정렬 결과와 공통 LLM
선택 절차가 결합된 파이프라인의 결과이다.

공개 코드 CE의 F1은 48.72점으로 수정 CE보다 9.85점 낮았다.
이 차이는 목록 참조 인덱스 수정에 따른 것으로, Jev 교체의 3.19점과
별개의 결과다. 문제를 확인한 범위는 공개 코드 스냅샷이며,
원 논문의 실제 실행에도 같은 문제가 있었는지는 확인되지 않았다.
